% ===================================================================
%  BẢNG Số 6 — Bên trong ngôi nhà, đồ đạc, việc ăn uống, việc thắp
%              sáng.
%
%  Traduction, faite en 2026, en vietnamien standard d'aujourd'hui,
%  sur l'ido de texto/io. Regles de la colonne : voir 10-bang-01.tex.
%  Cinq pieces, cinq numerotations : les cles portent c1 a c5.
%
%  LE TABLEAU LE PLUS LONG DU LIVRET, ET CELUI OU LE VIETNAMIEN
%  COMPTE. Deux cent cinquante objets defilent en cinq pieces, et
%  presque chacun demande un CLASSIFICATEUR — le petit mot qui dit de
%  quelle espece est la chose avant de la nommer :
%
%    cái   les objets sans vie      cái bàn la table, cái ghế la chaise
%    chiếc les objets uniques,      chiếc gương le miroir, chiếc khăn
%          souvent longs ou plats   la nappe
%    tấm   ce qui est plat et large tấm thảm le tapis, tấm rèm le rideau
%    bộ    ce qui va par jeu        bộ đồ ăn le couvert
%    con   ce qui bouge ou coupe    con dao le couteau
%
%  « Con dao » — le couteau prend le classificateur des etres animes,
%  comme le poisson et le chat — est la seule irregularite du tableau,
%  et elle est ancienne : ce qui tranche etait tenu pour vivant.
%
%  LA CUISINE EST VIETNAMIENNE, LE SALON EST FRANCAIS, ET C'EST LA
%  MEME LIGNE QU'AU TABLEAU 1. Nồi la marmite, chảo la poele, rổ le
%  panier, chổi le balai, thớt le billot, muôi la louche, dao le
%  couteau : rien n'a eu besoin d'etre emprunte, ces objets etaient
%  la. Mais le salon fait entrer d'un coup pi-a-nô le piano, đi-văng
%  le divan, an-bom l'album, ga-lăng la galerie, và-li la valise —
%  et « bình phong », le paravent, qui est chinois comme l'objet.
%
%  LA SALLE DE BAIN N'A PAS DE NOM ANCIEN, et il faut le dire : la
%  maison vietnamienne de 1926 n'avait pas de piece pour se baigner.
%  « Phòng tắm » est un compose transparent — piece-baigner — forge
%  sur le modele des autres pieces, et c'est bien ainsi qu'il est
%  entre dans la langue. Le chauffe-eau (2) et la douche (7) sont
%  restes francais, parce qu'ils sont arrives avec la piece.
% ===================================================================

%%K t06-apar-1 apar
\VUsaut{6.0mm}
\VUtitre{15.0pt}{60}{BẢNG Số 6}
\VUsaut{1.4mm}
\VUtitre{11.4pt}{0}{Bên trong ngôi nhà, đồ đạc, việc ăn uống,}
\VUsaut{0.4mm}
\VUtitre{11.4pt}{0}{việc thắp sáng.}
\VUsaut{2.2mm}

%%K t06-apar-2 apar
Ngôi nhà đã xong. Bác của Gioan đã dọn vào chỗ ở mới, mà ta đã biết cách bố trí. Bây giờ
ta hãy xem bác đã bày biện đồ đạc thế nào. Trước hết ta thăm :

%%K t06-tit-1 sub
\VUpk{10.2pt}{40}{Phòng ngủ.}
\VUsaut{3.0mm}

%%K t06-c1-01-1 p
1. --- Chúng ta đến không đúng lúc, vì chị hầu phòng đang bận dọn \VUgras{phòng}.

%%K t06-c1-02-1 p
2. --- Trên \VUgras{bàn rửa mặt} \textsuperscript{(1)}, đây đó là những đồ vật dùng để
rửa, để chải đầu, tóm lại là để làm vệ sinh. Đó là cái \VUgras{chậu}
\textsuperscript{(2)} và cái \VUgras{bình nước} \textsuperscript{(3)}, cái \VUgras{lược
chải đầu} \textsuperscript{(4)}, cái \VUgras{lược} \textsuperscript{(5)}, bánh
\VUgras{xà phòng} \textsuperscript{(6)} và miếng \VUgras{bọt biển}
\textsuperscript{(7)}, sau cùng là một \VUgras{lọ} \textsuperscript{(8)} nhỏ đựng nước
súc miệng.

%%K t06-c1-03-1 p
3. --- Dưới đất, trước bàn rửa mặt, kìa cái \VUgras{xô} \textsuperscript{(9)} để hứng
nước bẩn.

%%K t06-c1-04-1 p
4. --- Ở \VUgras{phòng ngoài} \textsuperscript{(10)}, ta thấy mấy cái \VUgras{móc áo}
\textsuperscript{(13)} và hai chiếc \VUgras{ô} \textsuperscript{(11)} trong \VUgras{giá
để ô} \textsuperscript{(12)}.

%%K t06-c1-05-1 p
5. --- Phía bên kia cửa là chiếc \VUgras{tủ có gương} \textsuperscript{(14)} đựng
\VUgras{đồ vải} \textsuperscript{(15)}.

%%K t06-c1-06-1 p
6. --- Ta hãy nhân lúc \VUgras{chị hầu phòng} \textsuperscript{(6)} dọn \VUgras{giường}
\textsuperscript{(17)} mà xem các bộ phận của nó. \VUgras{Khung giường}
\textsuperscript{(20)} đỡ một tấm \VUgras{đệm lò xo} \textsuperscript{(21)} tốt, trên đó
chị Iustina sắp đặt một tấm \VUgras{đệm len} \textsuperscript{(22)}. Rồi chị sẽ lấy trên
ghế hai tấm \VUgras{ga trải} \textsuperscript{(23)} và cái \VUgras{chăn}
\textsuperscript{(24)}. Sau đó chị đặt ở đầu giường chiếc \VUgras{gối dài}
\textsuperscript{(25)} và chiếc \VUgras{gối đầu} \textsuperscript{(26)} đang nằm cùng
\VUgras{chăn lông} \textsuperscript{(27)} trên \VUgras{tấm thảm chân giường}
\textsuperscript{(32)}. Chị dùng chổi lông để quét bụi \VUgras{rèm}
\textsuperscript{(19)} và \VUgras{trướng giường} \textsuperscript{(18)}, và thế là xong
một phần công việc.

%%K t06-c1-07-1 p
7. --- Rồi chị sẽ phải dọn lại chút ít \VUgras{bàn ngủ} \textsuperscript{(28)}, lên dây
\VUgras{đồng hồ báo thức} \textsuperscript{(29)}, sửa soạn \VUgras{đèn ngủ}
\textsuperscript{(30)}, lấy \VUgras{cây nến} \textsuperscript{(31)} ra để lau chân nến.

%%K t06-tit-2 sub
\VUsaut{5.0mm}
\VUpk{10.2pt}{40}{Giỏ đồ khâu và đồ dùng bên lò sưởi.}
\VUsaut{3.0mm}

%%K t06-c1-08-1 p
8. --- \VUgras{Giỏ đồ khâu} \textsuperscript{(34)} bên cạnh chiếc \VUgras{ghế đẩu}
\textsuperscript{(33)} cũng rất cần được dọn. Chị sẽ phải gấp tấm \VUgras{đồ thêu}
\textsuperscript{(35)}, cất vào \VUgras{bàn khâu} \textsuperscript{(38)} cái
\VUgras{đê}, cuộn \VUgras{chỉ} \textsuperscript{(36)}, những cây \VUgras{kim}
\textsuperscript{(37)} và cái \VUgras{kéo} \textsuperscript{(39)}, rồi khóa nắp bàn bằng
cái \VUgras{chìa} \textsuperscript{(40)}.

%%K t06-c1-09-1 p
9. --- Trên chiếc \VUgras{bàn một chân} \textsuperscript{(41)} ở giữa phòng, chị Iustina
sẽ thay nước trong \VUgras{bình} \textsuperscript{(42)}. Chị sẽ bỏ \VUgras{đường} vào
\VUgras{hộp đường} \textsuperscript{(43)}, lau cái \VUgras{kẹp đường}
\textsuperscript{(44)} và cất cái \VUgras{bàn chải quần áo} \textsuperscript{(46)}.

%%K t06-c1-10-1 p
10. --- Phía phải của phòng sẽ đỡ việc hơn cho chị, vì chiếc \VUgras{ghế dài}
\textsuperscript{(47)} và cái \VUgras{gối} \textsuperscript{(48)} của nó đã đúng chỗ, và
vì trời không lạnh nên không có gì bị xê dịch trong đám đồ dùng của \VUgras{lò sưởi}
\textsuperscript{(49)}. Cái \VUgras{xẻng} \textsuperscript{(50)}, cái \VUgras{kẹp than}
\textsuperscript{(51)}, đôi \VUgras{giá đỡ củi} \textsuperscript{(55)} và \VUgras{tấm
chắn tro} \textsuperscript{(56)} đều sạch sẽ; \VUgras{tấm chắn lửa}
\textsuperscript{(54)} không bị dời đi và \VUgras{thùng củi} \textsuperscript{(52)} đầy
ắp \VUgras{củi} \textsuperscript{(53)}.

%%K t06-noto-1 noto
\VUnotes{143.2mm}{%
(*) Babo = đứa bé; tiếng Anh \textit{baby}, tiếng Pháp \textit{bébé}, tiếng Ý
\textit{bambino}.%
}

%%K t06-noto-2 noto
\VUnotes{143.2mm}{%
(*) Lambrequino = tiếng Pháp \textit{lambrequin}, tiếng Tây Ban Nha
\textit{lambrequines}, tiếng Bồ Đào Nha \textit{lambrequins}, và cả tiếng Đức lẫn tiếng
Anh \textit{lambrequins}; nghĩa là tiếng Đức, tiếng Anh, tiếng Pháp, tiếng Bồ Đào Nha và
tiếng Tây Ban Nha đều như nhau.%
}

%%K t06-c1-11-1 p
11. --- Tuy vậy chị sẽ phải quét bụi chiếc \VUgras{đồng hồ quả lắc}
\textsuperscript{(57)}, những cây \VUgras{giá nến} \textsuperscript{(58)} và các
\VUgras{khay đựng đồ lặt vặt} \textsuperscript{(59)}, rồi lau chiếc \VUgras{gương}
\textsuperscript{(60)}.

%%K t06-c1-12-1 p
12. --- Còn \VUgras{cái nôi} \textsuperscript{(61)} của em bé (của babo (*)) thì đã sẵn
sàng.

%%K t06-tit-3 sub
\VUsaut{5.0mm}
\VUpk{10.2pt}{40}{Phòng khách.}
\VUsaut{3.0mm}

%%K t06-c2-01-1 p
1. --- Bây giờ là phòng khách. Đó là một \VUgras{căn phòng} rất đẹp. Lò sưởi, ở góc
phải, được trang trí bằng một \VUgras{bức tượng nhỏ} \textsuperscript{(1)} tinh xảo và
hai chiếc \VUgras{bình} \textsuperscript{(3)} đẹp, và phủ một \VUgras{diềm vải}
\textsuperscript{(2)} (*) bằng nhung.

%%K t06-c2-02-1 p
2. --- Để tia lửa từ lò không bén vào thảm, người ta đặt trước lò một \VUgras{tấm chắn
tia lửa} \textsuperscript{(4)} bằng đồng.

%%K t06-c2-03-1 p
3. --- Ngay cạnh đó, một chiếc \VUgras{bàn viết} \textsuperscript{(5)} thanh nhã dành
cho phụ nữ đang mở. Chính ở đó \VUgras{bà chủ nhà} \textsuperscript{(23)} viết thư.

%%K t06-c2-04-1 p
4. --- Cửa sổ có \VUgras{rèm che kính} \textsuperscript{(6)} được vén lên bằng
\VUgras{dây buộc rèm} \textsuperscript{(8)} móc vào \VUgras{núm gắn tường}
\textsuperscript{(7)}, và những tấm rèm vải lớn phía trên có \VUgras{diềm}
\textsuperscript{(9)}.

%%K t06-c2-05-1 p
5. --- Trong hốc cửa sổ, một \VUgras{chậu cảnh} \textsuperscript{(11)} đựng một
\VUgras{cây} \textsuperscript{(10)} xanh, một cây cọ tuyệt đẹp mà lá vươn gần tới
\VUgras{đèn dầu} \textsuperscript{(12)}.

%%K t06-c2-06-1 p
6. --- \VUgras{Chao đèn} \textsuperscript{(13)} của cây đèn do Maria sắp đặt, cô
\VUgras{thiếu nữ} \textsuperscript{(14)} duyên dáng đang ngồi bên đàn dương cầm
\textsuperscript{(15)}. Ngồi thẳng lưng trên chiếc \VUgras{ghế đẩu}
\textsuperscript{(16)}, cô tập một bản nhạc. Cô rất giỏi và cẩn thận, như \VUgras{tủ
đựng bản nhạc} \textsuperscript{(17)} ngăn nắp hoàn hảo của cô chứng tỏ.

%%K t06-tit-4 sub
\VUsaut{5.0mm}
\VUpk{10.2pt}{40}{Đứa trẻ nghịch ngợm.}
\VUsaut{3.0mm}

%%K t06-c2-07-1 p
7. --- Em trai cô, Phaolô, tìm một cuốn \VUgras{sách} \textsuperscript{(19)} trong
\VUgras{tủ sách} \textsuperscript{(18)}. Em là một \VUgras{cậu bé}
\textsuperscript{(20)} hiếu động và khó chịu. Bám vào tủ sách để với cuốn sách ở quá
cao, em có nguy cơ làm rơi \VUgras{bức tượng bán thân} \textsuperscript{(21)} ở phía
trên.

%%K t06-c2-08-1 p
8. --- Em làm điều dại dột ấy nhân lúc mẹ em đang mải trò chuyện với một người bạn, một
\VUgras{bà} \textsuperscript{(22)} đến chơi nhà mấy hôm. Người mẹ ngồi trên chiếc
\VUgras{đi-văng} \textsuperscript{(24)}, cạnh chiếc \VUgras{ghế bành}
\textsuperscript{(25)}, còn bạn bà thì ngồi trên chiếc \VUgras{ghế}
\textsuperscript{(27)} mà ta chỉ thấy cái \VUgras{lưng tựa} \textsuperscript{(26)}.

%%K t06-c2-09-1 p
9. --- Chiếc \VUgras{khung ảnh} \textsuperscript{(30)} lớn treo trên tường chứa
\VUgras{bức chân dung} \textsuperscript{(29)} của một người bà con. Trong cuốn
\VUgras{an-bom} \textsuperscript{(32)} đặt trên bàn có tập hợp \VUgras{ảnh chụp}
\textsuperscript{(33)} của những người khác trong gia đình. Bọn trẻ vừa mới giở nó ra,
vì mấy chiếc \VUgras{ghế} \textsuperscript{(31)} còn quây quanh bàn.

%%K t06-c2-10-1 p
10. --- Chiếc \VUgras{bình sứ Trung Hoa} \textsuperscript{(34)} đặt gần con \VUgras{dao
rọc giấy} \textsuperscript{(35)} là quà tặng Maria nhân ngày lễ của cô, và tấm
\VUgras{bình phong} \textsuperscript{(36)} trang trí góc phòng do bác cô mang từ Trung
Hoa về.

%%K t06-c2-11-1 p
11. --- Được những \VUgras{bức tranh} \textsuperscript{(37)} tuyệt đẹp treo trên
\VUgras{vách} \textsuperscript{(42)} làm cho vui mắt, được chiếc \VUgras{đèn chùm}
\textsuperscript{(38)} với những \VUgras{bóng đèn điện} \textsuperscript{(39)} tỏa sáng
vui vẻ vào buổi tối, phòng khách ấy trông thật dễ chịu. \VUgras{Tấm thảm}
\textsuperscript{(40)} êm trải trên sàn gỗ và \VUgras{chuông điện}
\textsuperscript{(41)} rất tiện lợi làm cho nó thêm đầy đủ tiện nghi.

%%K t06-tit-5 sub
\VUsaut{3.6mm}
\VUpk{10.2pt}{40}{Phòng ăn.}
\VUsaut{3.0mm}

%%K t06-c3-01-1 p
1. --- Chuông báo bữa tối đã reo. Cả gia đình, trừ Maria, quây quần quanh bàn. Phòng ăn
được sưởi ấm dễ chịu nhờ chiếc \VUgras{lò sưởi} \textsuperscript{(1)} bằng sành rất tốt,
có \VUgras{ống khói} \textsuperscript{(2)} nhỏ.

%%K t06-c3-02-1 p
2. --- Căn phòng này không có nhiều đồ đạc. Dọc tường, phía trên chiếc \VUgras{ghế
kê chân} \textsuperscript{(4)}, treo một \VUgras{bể nước nhỏ} \textsuperscript{(3)}
trang trí. Ngay cạnh là chiếc \VUgras{tủ bày} \textsuperscript{(5)}, nơi người ta đặt
các món nguội, món tráng miệng và những đồ dùng chưa cần đến ngay.

%%K t06-c3-03-1 p
3. --- Như vậy, trên ngăn trên ta thấy một con \VUgras{gà quay} \textsuperscript{(7)},
một con \VUgras{cá} \textsuperscript{(8)} và một chiếc \VUgras{cốc lớn}
\textsuperscript{(10)} đựng \VUgras{trái cây} \textsuperscript{(9)}. Ở dưới là
\VUgras{phô mai} \textsuperscript{(11)}, \VUgras{bánh mì} \textsuperscript{(12)} và
\VUgras{bộ đựng dầu giấm} \textsuperscript{(13)}, chứa mọi thứ cần để trộn xà lách :
dầu, giấm, \VUgras{muối} \textsuperscript{(14)} và \VUgras{hạt tiêu}
\textsuperscript{(15)}. Sau cùng, trên ngăn dưới là một chiếc \VUgras{bánh ngọt}
\textsuperscript{(17)} bên cạnh \VUgras{lọ mù tạt} \textsuperscript{(16)}.

%%K t06-c3-04-1 p
4. --- Chiếc đồng hồ của phòng ăn, gọi là \VUgras{đồng hồ treo tường}
\textsuperscript{(20)}, có hình dáng rất riêng. Nó được gắn trong một khung gỗ, khung ấy
còn chứa cả \VUgras{phong vũ biểu} \textsuperscript{(19)} và \VUgras{nhiệt kế}
\textsuperscript{(18)}.

%%K t06-tit-6 sub
\VUsaut{4.6mm}
\VUpk{10.2pt}{40}{Bữa tối được dọn.}
\VUsaut{3.0mm}

%%K t06-c3-05-1 p
5. --- Dọc tất cả các bức tường của phòng chạy một đường \VUgras{ván ốp}
\textsuperscript{(21)} bằng gỗ sồi đánh sáp. Một \VUgras{tấm rèm cửa}
\textsuperscript{(22)} bằng vải che kín cánh cửa thông ra hiên. Chính ở đó gia đình sẽ
uống cà phê sau bữa tối. \VUgras{Tách} \textsuperscript{(24)} và \VUgras{đĩa lót}
\textsuperscript{(25)} đã sẵn sàng trên chiếc \VUgras{tủ buýp phê}
\textsuperscript{(23)}.

%%K t06-c3-06-1 p
6. --- Phía trên bàn, gắn vào trần là một \VUgras{đèn treo} \textsuperscript{(26)} mà
ngọn đèn rất sáng của nó soi cả phòng ăn vào buổi tối.

%%K t06-c3-07-1 p
7. --- Bây giờ ta hãy xem chính chiếc \VUgras{bàn ăn} \textsuperscript{(27)}. Khi gia
đình bước vào, bàn đã dọn xong. Trên \VUgras{khăn trải bàn} \textsuperscript{(28)},
trước chỗ mỗi người, người ta đã đặt một cái \VUgras{đĩa} \textsuperscript{(33)}, một
chiếc \VUgras{khăn ăn} \textsuperscript{(29)}, một cái \VUgras{thìa}
\textsuperscript{(34)}, một cái \VUgras{nĩa} \textsuperscript{(35)}, một con
\VUgras{dao} \textsuperscript{(36)} và một cái \VUgras{ly} \textsuperscript{(32)}. Cũng
có những \VUgras{chai} \textsuperscript{(30)} đựng rượu vang và một \VUgras{bình}
\textsuperscript{(31)} đựng nước.

%%K t06-c3-08-1 p
8. --- \VUgras{Người giúp việc} \textsuperscript{(39)} trước hết mang lên \VUgras{liễn
xúp} \textsuperscript{(37)}, trong đó xúp còn bốc khói. Bây giờ anh dọn \VUgras{món}
\textsuperscript{(38)} thứ nhất, món thịt. Rồi anh sẽ dọn những món ta đã kể; sau cùng,
sau \VUgras{món tráng miệng}, anh sẽ nhân lúc chủ nhà sang hiên để dọn bát đĩa và sắp
xếp lại phòng ăn.

%%K t06-c3-08-2 p
\textit{Chú thích.} --- \textit{Những từ thông dụng liên quan đến việc ăn uống ở đây chỉ
được nêu rất vắn tắt. Danh sách sẽ được bổ sung trong hai bảng « Khách sạn »
(số 13) và « Chợ » (số 15).}

%%K t06-tit-7 sub
\VUsaut{4.6mm}
\VUpk{10.2pt}{40}{Nhà bếp.}
\VUsaut{3.0mm}

%%K t06-c4-01-1 p
1. --- Trong bếp, mà ta nhìn qua cánh cửa hé mở, ta thấy một cảnh bừa bộn nên thơ. Trước
\VUgras{bếp lửa} \textsuperscript{(1)}, nơi \VUgras{lửa} \textsuperscript{(3)} nổ lách
tách, đặt một cái \VUgras{quay thịt} \textsuperscript{(5)}. Một tảng \VUgras{thịt quay}
\textsuperscript{(7)} ngon lành đang \VUgras{xiên} \textsuperscript{(6)} và sắp chín.

%%K t06-c4-02-1 p
2. --- Cái \VUgras{bễ thổi} \textsuperscript{(8)} và cái \VUgras{xẻng than}
\textsuperscript{(9)} vừa dùng xong còn nằm dưới đất, cùng với những \VUgras{thanh củi}
\textsuperscript{(10)}.

%%K t06-c4-03-1 p
3. --- Mặt trên lò sưởi đã gọn gàng : cái \VUgras{đèn lồng} \textsuperscript{(11)}, cái
\VUgras{hộp diêm} \textsuperscript{(13)} đựng những que \VUgras{diêm}
\textsuperscript{(12)}, cái \VUgras{chân nến} \textsuperscript{(14)} và cái \VUgras{đế
nến} \textsuperscript{(15)} với cây \VUgras{nến} \textsuperscript{(16)} đều đúng chỗ.
Nhưng chiếc \VUgras{xô than} \textsuperscript{(18)} lớn, đầy \VUgras{than đá}
\textsuperscript{(17)} mà \VUgras{chị bếp} \textsuperscript{(23)} vừa xuống tầng hầm lấy
lên, còn nằm giữa bếp.

%%K t06-c4-04-1 p
4. --- Quả thật, người đàn bà đáng thương phải vội vàng để bữa trưa kịp giờ. Lúc này chị
đang ở gần \VUgras{bếp lò} \textsuperscript{(19)} và khuấy một món \VUgras{nước xốt}
\textsuperscript{(24)} không được để cháy quá.

%%K t06-c4-05-1 p
5. --- Trong \VUgras{lò nướng} \textsuperscript{(20)} có một chiếc bánh chị làm cho bữa
xế của bọn trẻ. Phía trên bếp lò treo cái \VUgras{muôi} \textsuperscript{(21)} và cái
\VUgras{muôi thủng} \textsuperscript{(22)}.

%%K t06-tit-8 sub
\VUsaut{4.0mm}
\VUpk{10.2pt}{40}{Bát đĩa xoong nồi.}
\VUsaut{3.0mm}

%%K t06-c4-06-1 p
6. --- \VUgras{Bộ đồ nấu bếp} \textsuperscript{(25)} treo ở cuối phòng rất sạch sẽ.
Những chiếc \VUgras{xoong} \textsuperscript{(26)} đồng đẹp đẽ, cái \VUgras{bình sữa}
\textsuperscript{(27)}, cái \VUgras{rây} \textsuperscript{(28)} và cái \VUgras{bàn nạo}
\textsuperscript{(29)} đã được lau chùi kỹ lưỡng.

%%K t06-c4-07-1 p
7. --- Cái \VUgras{hũ muối} \textsuperscript{(30)} đặt trên chiếc tủ nhỏ, cạnh cái
\VUgras{ấm trà} \textsuperscript{(31)} và cái \VUgras{vò} \textsuperscript{(32)} đựng
dầu. Cái \VUgras{ngăn kéo} \textsuperscript{(33)} chắc là đựng bộ đồ ăn và dao bếp.

%%K t06-c4-08-1 p
8. --- Cái \VUgras{chổi} \textsuperscript{(34)} và cái \VUgras{chổi lông}
\textsuperscript{(35)} ở một góc. Chị bếp vừa dùng cái \VUgras{thớt}
\textsuperscript{(36)}; chị để nó gần cửa sổ.

%%K t06-c4-09-1 p
9. --- Một chiếc \VUgras{khăn lau tay} \textsuperscript{(37)} treo trên tường, cạnh cái
\VUgras{phễu} \textsuperscript{(38)}. Chính ở góc bếp này người ta để cái \VUgras{vỉ
nướng} \textsuperscript{(39)}, con \VUgras{dao bầm} \textsuperscript{(40)} và tấm
\VUgras{thớt bầm} \textsuperscript{(41)}. Rồi phía trên dao bầm, trên một tấm \VUgras{ván}
\textsuperscript{(42)}, có những cái \VUgras{hũ} \textsuperscript{(43)}, cái \VUgras{nồi
hầm} \textsuperscript{(44)} với cái \VUgras{vung} \textsuperscript{(45)} và cái
\VUgras{nồi to} \textsuperscript{(46)}. Cái \VUgras{chảo} \textsuperscript{(47)} treo
ngay cạnh.

%%K t06-c4-10-1 p
10. --- Chỉ có cái \VUgras{vòi nước} \textsuperscript{(48)} bằng đồng là sáng lên trong
góc này. Cái \VUgras{bồn rửa} \textsuperscript{(49)}, trên đó đặt cái \VUgras{vại}
\textsuperscript{(50)} to, chắc rất tiện lợi với chiếc tủ nhỏ, nơi người ta cất
\VUgras{thùng giấm} \textsuperscript{(51)} có \VUgras{vòi} \textsuperscript{(52)}, cái
\VUgras{thùng rác} \textsuperscript{(53)} và \VUgras{bát đĩa bẩn} \textsuperscript{(54)}.

%%K t06-tit-9 sub
\VUsaut{4.0mm}
\VUpk{10.2pt}{40}{Những đồ khác để trong bếp.}
\VUsaut{3.0mm}

%%K t06-c4-11-1 p
11. --- Cái \VUgras{chậu lớn} \textsuperscript{(55)} dùng để rửa \VUgras{rau}
\textsuperscript{(58)} và cái \VUgras{nồi gang} \textsuperscript{(56)} đặt trên
\VUgras{nền gạch} \textsuperscript{(57)} chắc cũng có chỗ của chúng trong chiếc tủ ấy.

%%K t06-c4-12-1 p
12. --- Chị bếp chắc chắn không thiếu lửa, vì kìa, trên góc bàn, còn có một cái
\VUgras{bếp ga} \textsuperscript{(59)} đang đun nước trong một chiếc xoong nhỏ.
\VUgras{Hơi nước} \textsuperscript{(60)} thoát ra cho ta biết nước chắc đang sôi. Nước ấy
hẳn là để pha cà phê, vì kìa, ở đầu kia của bàn, là cái \VUgras{ấm cà phê}
\textsuperscript{(64)} và cái \VUgras{cối xay cà phê} \textsuperscript{(63)}.

%%K t06-c4-13-1 p
13. --- Bếp ga nối với \VUgras{vòi ga} \textsuperscript{(62)} bằng một \VUgras{ống cao
su} \textsuperscript{(61)} dài.

%%K t06-c4-14-1 p
14. --- \VUgras{Chị giúp việc theo ngày} \textsuperscript{(66)} đến giúp chị bếp đang
làm rau cho bữa tối. Khi rau đã nhặt và rửa xong, chị sẽ bỏ vào cái \VUgras{chậu men}
\textsuperscript{(65)} chờ đến giờ nấu.

%%K t06-tit-10 sub
\VUsaut{4.0mm}
{\centering\fontsize{10.2pt}{10.2pt}\selectfont\textls[40]{\scshape Phòng tắm.} \textit{(Buồng tắm.)}\par}
\VUsaut{3.0mm}

%%K t06-c5-01-1 p
1. --- Chúng ta chỉ còn một chút tò mò nữa để biết hết những phòng chính của căn nhà :
xem cách bố trí phòng tắm. Việc ấy không lâu, vì phòng không rộng, và ta sẽ nhanh chóng
biết rằng \VUgras{bồn tắm} \textsuperscript{(1)} có một \VUgras{bình đun nước}
\textsuperscript{(2)} tốt, rằng \VUgras{đĩa đựng xà phòng} \textsuperscript{(3)} nằm vừa
tầm tay người tắm, và rằng \VUgras{giá treo khăn} \textsuperscript{(5)} chắc chắn giúp
\VUgras{khăn tắm} \textsuperscript{(4)} mau khô.

%%K t06-c5-02-1 p
2. --- Căn phòng này có vách ốp \VUgras{gạch men} \textsuperscript{(6)}, nhờ đó người ta
lắp được một \VUgras{vòi sen} \textsuperscript{(7)} mà không sợ nước bắn ra khi dùng làm
hỏng tường. Một cái \VUgras{chậu} \textsuperscript{(8)} nhỏ bằng kẽm, dùng để ngâm chân,
làm cho đồ đạc thêm đầy đủ.
\VUsaut{6.0mm}
\VUfilet{22mm}
